\begin{problem}
    Un estanque de agua tiene forma de prisma recto de base rectangular, cuya base mide \SI{4}{m} de ancho y \SI{3}{m} de largo. Este se debe llenar hasta \SI{2}{m} de altura con agua proveniente de un ducto.

    Para calcular el caudal del agua que sale de un ducto, se puede utilizar la siguiente expresión:
    \[
    \frac{V}{t}
    \]
    tal que $V$ corresponde al volumen de agua que sale del ducto durante $t$ unidades de tiempo.

    ¿Cuál de las siguientes expresiones corresponde al caudal de agua que debe salir del ducto, medido en \si{m^3/h}, para que el estanque se llene hasta \SI{2}{m} de altura en \SI{10}{h}?

    \begin{enumerate}[label=\Alph*.]
        \item $\dfrac{4 \cdot 3 \cdot 2}{10}$
        \item $\dfrac{4 + 3 + 2}{10}$
        \item $\dfrac{2(4 + 3)}{10}$
        \item $\dfrac{2(4 \cdot 3 + 2 \cdot 4 + 2 \cdot 3)}{10}$
    \end{enumerate}
\end{problem}
